\documentclass{report}
\usepackage{amsfonts, amsmath}
\newcommand\R{{\mathbb R}}
\newcommand\Q{{\mathbb Q}}
\newcommand\N{{\mathbb N}}
\newcommand\ve{\varepsilon}
\newcommand\Co{{\mathcal C}}
\pagestyle{empty}
\begin{document}
\large
\centerline{\Huge Fisica Matematica II}
\renewcommand{\thefootnote}{ } 
\footnote{\sl Avete 2 ore di tempo. Ogni esercizio
vale dieci punti. La sufficienza si ottiene con un punteggio $\geq$
18. Solo le risposte chiaramente giustificate saranno prese in
considerazione.}
\renewcommand{\thefootnote}{\arabic{footnote}} 
\vskip.1cm
\centerline{\Large Esonero 30-04-2003}
\vskip1cm
\begin{enumerate}
\item Si consideri il seguente problema di Cauchy
\[
(1-y^2)u_y+u_x=u
\]\[
u(x,0)=x.
\]
Discutere esistenza, unicit\`a e dominio delle soluzioni.
\item Si classifichino tutte le funzioni armoniche su $\R^2$ della forma
$f(x)g(y)$.

\item Si consideri il seguente problema di Dirichlet sul dominio
$\Omega=\{(x,y)\in\R^2\;|\; 1< 4x^2+y^2< 4\}$.
\[
u_{xx}+4u_{yy}=0\]\[  
u(x,y)=f(x)\quad\text{per }4x^2+y^2=1\]\[
u(x,y)=g(y)\quad\text{per }4x^2+y^2=4.
\]
Si discuta l'esitenza e l'unicit\`a. Si calcoli
\[
\int_{\Omega} u_x,
\]
in funzione di $f, g$.
Infine, si trovi la soluzione per $f(x)=0$, $g(y)=1$.
\item Si consideri il cambio di variabili $\Psi:\R^n\to\R^n$ definito
da $\Psi(x)=Ax+b$. Si calssifichino tutti gli $A, b$ per cui $\Psi$
manda funzioni armoniche in funzioni armoniche.
\end{enumerate}
\end{document}
